\subsection{Jacobian-Based Contraction on an Invariant Domain}
\label{sec:implemented-domain-contraction}

\subsubsection{Pipeline Overview}
The starting point is a randomly sampled equation that uses past observations
to generate the next value. We adjust its coefficients so that it can be
repeated safely within a chosen range. The pipeline is:
\begin{enumerate}
\item \textbf{Sample the equation.} Generate an AST with lagged inputs and
coefficients, preserving an explicit formula for the mechanism.
\item \textbf{Choose the allowed range and noise.} Specify where all initial
values may lie and the largest innovation that may be added at each step.
\item \textbf{Bound the output and its sensitivity.} Calculate conservative
limits on the equation's magnitude and on how changes in its historical
inputs affect the next state, including the copied lag values.
\item \textbf{Scale and verify the equation.} Reduce its overall strength
until the generated values remain in the allowed range and differences
between histories shrink under shared future noise.
\item \textbf{Generate trajectories.} Draw initial values inside that range
and execute the adjusted equation with bounded random innovations.
\item \textbf{Check the implementation empirically.} Simulate several
trajectories and check that their values and measured sensitivities respect
the calculated limits.
\end{enumerate}
Steps 3--4 establish the sufficient mathematical conditions; Step 6 tests
whether the implementation behaves consistently with them. The following
steps explain both the calculation and its purpose.

\subsubsection{Steps 1--2: Sample the Equation and Define Its Allowed Inputs}
This method controls the feedback created when generated values become inputs
to later steps. We want two properties: every observation stays inside a
chosen interval, and differences caused by the initial history shrink over
time. To examine the second property, run the same equation from two different
histories and give both runs exactly the same sampled noise value at each
future step. Their remaining difference then reflects the effect of the
initial histories. Fresh noise can keep both series fluctuating.

The sampler in \path{src/tsfmxai/generators/contraction.py} generates a
sum/product AST and attaches coefficients to its lag leaves. We choose
an allowed interval $[-r_D,r_D]$ and a noise bound $b_\varepsilon<r_D$.
Every initial lag value must belong to that interval. These choices define
where the equation will be checked and reserve room for the innovations.

\paragraph{Store the values needed for the next update.}
Consider the two-lag example
\begin{equation}
 x_{t+1}=0.2x_t+0.1x_tx_{t-1}+\varepsilon_{t+1},
 \qquad F_{\mathcal T}(s_0,s_1)=0.2s_0+0.1s_0s_1.
 \label{eq:jacobian-explanatory-recurrence}
\end{equation}
The function $F_{\mathcal T}$ is the deterministic part of the equation;
$\varepsilon_{t+1}$ is the random innovation added at that step.
The next value requires both $x_t$ and $x_{t-1}$, so store them together as
$\mathbf{s}_t=(x_t,x_{t-1})^\top$. This vector is called the \emph{state}:
it contains the complete history needed for one update. Its first coordinate,
$s_0$, is the newest value; $s_1$ is the preceding value.

After generating $x_{t+1}$, the new state is $(x_{t+1},x_t)^\top$:
insert the new value, move the previous newest value one position back, and
discard the oldest value. When the largest lag is $L$, the state stores the
$L$ most recent observations, and this complete update is
\begin{equation}
 \mathcal F_\varepsilon(\mathbf{s})=
 \bigl(F_{\mathcal T}(\mathbf{s})+\varepsilon,
       s_0,\ldots,s_{L-2}\bigr)^\top.
 \label{eq:implemented-state-transition}
\end{equation}
Thus $\mathcal F_\varepsilon$ describes both the equation and its memory
update. Stacking lagged observations this way is the \emph{augmented-state}
representation discussed by \citet{laraba2016setinvariance}.

\subsubsection{Step 3: Bound the Output and the Sensitivity of the Whole State}
\paragraph{Use derivatives to measure sensitivity to that history.}
A partial derivative measures how an output changes when one input changes
slightly while the other inputs stay fixed. The \emph{Jacobian} collects
these derivatives for every output coordinate of the state update.
For the example above,
\begin{equation}
 J_{\mathcal F}(\mathbf{s})=
 \begin{pmatrix}
  \partial F_{\mathcal T}/\partial s_0 &
  \partial F_{\mathcal T}/\partial s_1\\
  1 & 0
 \end{pmatrix}
 =\begin{pmatrix}
   0.2+0.1s_1 & 0.1s_0\\
   1 & 0
  \end{pmatrix}.
 \label{eq:jacobian-explanatory-matrix}
\end{equation}
The first row measures the sensitivity of the newly generated value.
The second row describes memory: the previous newest value is copied with
multiplier one. The additive noise has derivative zero with respect to the
history. In the two-run comparison its value is shared, so it also cancels
when subtracting the two updates.

\paragraph{Evaluate the bounds over every allowed history.}
Choose a radius $r_D>0$ and require every stored value to lie in
$[-r_D,r_D]$. The set of all allowed histories is
$D=[-r_D,r_D]^L$, called a \emph{box}. With two lags and $r_D=1$, it contains
every pair $(s_0,s_1)$ whose coordinates are between $-1$ and $1$.
Write $\overline d_d$ for an upper bound on the magnitude of the derivative
with respect to coordinate $s_d$. On this box, the example's bounds are
$\overline d_0=0.3$ and $\overline d_1=0.1$.
They cover every allowed history, rather than only the values visited by
sampled trajectories.

Derivatives initially describe small changes. To control a finite difference
between two histories, consider the straight line joining them. Every point
on that line remains inside a box; this property is called \emph{convexity}.
The derivative bound therefore applies along the entire line and bounds the
total output difference. The argument requires a continuously differentiable
update: the derivatives exist and vary continuously. The sampled sums and
products satisfy this smoothness requirement.

\paragraph{Measure changes in the whole state, including copied lags.}
For $L\geq2$, each memory-copy row has one entry equal to one. The ordinary
infinity norm takes the largest absolute row sum, so the augmented Jacobian
has norm at least one even when the recurrence eventually forgets its
history. To establish strict contraction at each update, we use weights
that account for information moving into older memory positions while
retaining its influence on new observations. Specifically, we use fixed
positive weights $w_d=q_w^d$, with $q_w>1$, and measure a state
difference $\Delta\mathbf{s}=\mathbf{s}-\mathbf{s}'$ by
\begin{equation}
 \|\Delta\mathbf{s}\|_{\mathbf w,\infty}
 =\max_{0\leq d<L}\frac{|s_d-s'_d|}{w_d}.
 \label{eq:jacobian-weighted-distance}
\end{equation}
In words, divide each coordinate difference by its weight and take the largest
result. This is the \emph{weighted infinity norm}; the same weights are used
at every step. For two lags with $q_w=1.25$, an unchanged difference copied
from the newest to the older position is divided by 1.25, giving a factor
of 0.8. At the same time, the older lag's sensitivity in the new-value row is
multiplied by 1.25, so its effect on future observations remains included.

Combining the new-value row and the memory-shift rows gives, for $L\geq2$,
\begin{equation}
 \sup_{\mathbf{s}\in D}\|J_{\mathcal F}(\mathbf{s})\|_{\mathbf w,\infty}
 \leq\max\left\{\sum_{d=0}^{L-1}\overline d_d q_w^d,
                       q_w^{-1}\right\}<1.
 \label{eq:implemented-weighted-contraction}
\end{equation}
Here $\sup$ means the worst case over all histories in the box, and the matrix
norm measures the largest amplification of the state difference in
Eq.~\eqref{eq:jacobian-weighted-distance}. The sum bounds the new-value row;
$q_w^{-1}$ bounds each copy into the next history position. For $L=1$,
only the new-value row exists. Call a verified common upper bound $\gamma<1$.
While both runs remain in $D$, one update reduces the measured distance by
a factor at most $\gamma$. Repeating this comparison for $v$ updates gives
\begin{equation}
 \|\mathbf{s}_{\tau+v}-\mathbf{s}'_{\tau+v}\|_{\mathbf w,\infty}
 \leq\gamma^v
       \|\mathbf{s}_{\tau}-\mathbf{s}'_{\tau}\|_{\mathbf w,\infty}.
 \label{eq:jacobian-history-forgetting}
\end{equation}
The distance decreases geometrically. For example, $\gamma=0.8$ limits it
after ten updates to $0.8^{10}\approx0.107$ times its initial value.
This is a guaranteed upper bound on the remaining distance; particular
trajectories can lose their dependence on the initial history faster.
Each coordinate difference is at most its fixed weight multiplied by this
measured distance. As the distance tends to zero, every individual
lag-coordinate difference therefore tends to zero as well. The weighting
changes how we measure the histories while leaving their generated values
unchanged.

\paragraph{Bound the output to keep subsequent states inside the checked box.}
The bound above remains useful only while the states stay inside $D$.
We therefore also require
$\mathcal F_\varepsilon(D)\subseteq D$ for every allowed innovation: starting
inside the box must produce another state inside it. This property is called
\emph{invariance}. With $|\varepsilon_t|\leq b_\varepsilon$, it is ensured by
\begin{equation}
 \sup_{\mathbf{s}\in D}|F_{\mathcal T}(\mathbf{s})|
       +b_\varepsilon\leq r_D.
 \label{eq:jacobian-domain-containment}
\end{equation}
The equation reserves enough room for both the deterministic output and the
largest allowed noise value. The copied lag values already belong to the
interval. Starting with a history in $D$, the same reasoning applies again
at every step: all observations stay bounded, and the sensitivity estimate
remains applicable.

\paragraph{Connection to the published theory.}
\citet[Definition~6 and Theorems~15--16]{tran2019convergence} develop the
relationship between contraction and \emph{incremental stability}, the
decrease of differences between trajectories. The fixed weighted norm and
box argument above are the project's sufficient construction for this
property. The state stacking and invariant-set viewpoint follow the concepts
reviewed by \citet{laraba2016setinvariance}. The conclusions here concern
histories inside the specified box and the specified bounded innovations.

\paragraph{Calculate conservative bounds from the AST.}
Starting at the leaves, the code records a possible value interval and
absolute derivative bounds for each subtree. Addition adds the bounds;
multiplication combines the intervals and uses the product rule for the
derivatives. For example, $s_0s_1$ lies in $[-1,1]$ when both inputs do, and
its derivatives are $s_1$ and $s_0$, each bounded in absolute value by one.
These calculations cover the full box. They can overestimate the actual
range, particularly when repeated leaves refer to the same coordinate.

\subsubsection{Step 4: Scale the Equation and Verify Both Conditions}
The constructor multiplies the complete expression by a calculated factor.
Reducing this factor decreases both the
deterministic output bound and the new-value row's derivative bounds. The
factor must leave room for the noise in
Eq.~\eqref{eq:jacobian-domain-containment} and meet the target gain in
Eq.~\eqref{eq:implemented-weighted-contraction}. The copied-lag gain depends
on $q_w$, so this weight parameter is chosen to meet the target as well.
The final expression retains the multiplier, making the executed equation
available for inspection and ground-truth extraction.

\paragraph{Worked example using the implemented rules.}
For Eq.~\eqref{eq:jacobian-explanatory-recurrence} on $[-1,1]^2$, the
unscaled deterministic magnitude is at most 0.3 and the derivative bounds
are $(0.3,0.1)$. The comparison settings use $q_w=1.25$, target gain 0.9,
and uniform innovations with
$b_\varepsilon=\sqrt{3}\,0.05\approx0.0866$. This uniform law has standard
deviation 0.05, matching the Gaussian noise scale in the other families.
These numerical settings are implementation choices within the admissible
conditions: the target 0.9 leaves a gap of 0.1 below the contraction threshold,
and $q_w=1.25$ gives a memory-copy gain of 0.8, below that target. The safety
multiplier 0.95 applies a further 5\% reduction to the allowed drift scale.
Both checks already pass; the constructor still applies its 0.95 safety
multiplier. The executed drift is therefore
$0.19s_0+0.095s_0s_1$. Its derivative bounds are $(0.285,0.095)$, giving
\begin{equation}
\begin{gathered}
 \gamma=\max\{0.285+1.25(0.095),\,0.8\}=0.8,
 \\
 0.285+\sqrt{3}(0.05)\approx0.3716<1.
\end{gathered}
 \label{eq:jacobian-worked-checks}
\end{equation}
Here the memory-copy row determines the final gain of 0.8; scaling the drift
reduces only the new-value row. The first check controls forgetting of the
initial history; the second
keeps the generated values inside $[-1,1]$. This is a hand-worked project
example, checked against the constructor in the test suite.

\paragraph{Retain the verification record.}
After scaling, the constructor checks the adjusted equation and stores the
domain, noise bound, derivative bounds, and resulting gain. These quantities
describe the equation that will actually be simulated. The implementation
uses floating-point calculations and a safety margin when evaluating the
bounds.

\paragraph{Effect on the sampled mechanisms.}
The scale changes the coefficients of the sampled equation and therefore
the dynamics represented in the benchmark. Older lags receive larger
multipliers in the sensitivity bound: with $q_w=1.25$, offset $d=11$
(lag 12) contributes with weight $1.25^{11}\approx11.64$. This can require
stronger attenuation of the whole expression. With the innovation scale
held fixed, that attenuation can weaken dependence on the history and
increase the relative influence of noise. These possible effects motivate
the amplitude and nearly constant-tail diagnostics in
Section~\ref{sec:generator-diagnostics}. The bound is a sufficient condition;
an expression exceeding it may still satisfy a different stability criterion.

\subsubsection{Step 5: Draw Initial Histories and Generate Trajectories}
The current constructor handles one series and excludes an explicit time
terminal. Its initial history is sampled uniformly inside the box; subsequent
observations are generated directly. Uniform innovations supply the finite
noise bound required for containment in the box.
The mathematical checks permit any fixed initial history inside the box;
uniform sampling is the current implementation's choice. A larger allowed
range requires recalculating the bounds and the scale. The trajectory
generator also checks that each newly generated value remains finite and
inside the certified interval.

\paragraph{Stationary law and initialization.}
For each fixed sampled mechanism, the contraction and bounded invariant
domain above, together with i.i.d. innovations, satisfy the conditions of
\citet[Theorem~1]{diaconis1999iterated} for a unique stationary distribution
of the lag state. The state distribution approaches this law from an
arbitrary initial history inside the box. Strict stationarity from the
start requires drawing the initial lag state jointly from that stationary
law, independently of future innovations. Our uniform history sampling can
instead produce an initialization transient, during which the distribution
changes as the process approaches its stationary law. As discussed in
Section~\ref{sec:generator-properties}, containment and forgetting still
hold throughout that transient.

\subsubsection{Step 6: Test the Implementation by Simulation}
Monte Carlo simulations exercise the generator using several sampled
histories and innovation sequences. They check whether generated values
remain inside the chosen interval and whether local Jacobian values respect
the verified gain. This can reveal an error in the bound calculation, lag
ordering, or trajectory execution. The mathematical guarantee comes from
the full-domain checks in Steps 3--4; the simulations provide an additional
implementation check. The common diagnostics in
Section~\ref{sec:generator-diagnostics} then measure explosions, low amplitude,
and nearly constant tails under the declared evaluation protocol.
